\scanpage{16}
\diagram{\begin{tikzcd}[ampersand replacement=\&,row sep=small]
X\arrow[d]\\ Y\ \deleted{\unsure{G|6}}
\end{tikzcd}}
\[
\mathrm{Pic}(X/Y)
\]

$\mathrm{Pic}_{\mathrm{sp}\deleted{\unsure{é}}}(X/Y)$ : classes à \unsure{isom|isom.} près des faisceaux inversibles $\underline L$, section marquée, telles que \deleted{\unsure{les} \unsure{div} soit} \deleted{(i)} l'on ait
\[
H^{i}_{f}\deleted{\unsure{(X,}}(\underline L\deleted{\unsure{(n))}})=0
\quad\text{pour }i>0,\ \deleted{\unsure{n\gg0}}
\]
donc $\underline H_f^0(\underline L)$ est loc libre de rang $\chi_{X/Y}(\underline L)$.
\deleted{$(ii)\quad\unsure{f_{*}}\ \unsure{\underline H^0_f}\ \chi_{X/Y}(\underline L)$}

% Wavy cancellation through the lower right-hand block.
\deleted{\unsure{on prend} $>1$, i.e l'\unsure{hom|hyp.} : \unsure{déduit|de div} de la section marquée $H^0_f(\underline L)\longrightarrow$}

\[
\underline H^0_f(\underline{\mathrm O}_X)=\underline{\mathrm O}_Y
\]
\deleted{\unsure{et}} ``universellement''

\unsure{une|ou} section marquée :

% This diagram is upside down relative to the definition.
\diagram{\begin{tikzcd}[ampersand replacement=\&,column sep=small,row sep=small]
\deleted{\unsure{P}}\&{}\&{}\\
\mathcal A\arrow[dd]\&\mathcal D\arrow[l]\&
\mathcal D\times_{\mathcal A}\mathcal D\arrow[l,shift left=.5ex]\arrow[l,shift right=.5ex]\\
{}\&{}\&\mathcal D\times_Y\mathcal D\arrow[u,phantom,"\cap"]\\
Y\& P\arrow[uu]\&
P\times_Y\deleted{\unsure{P}}\ \unsure{D_1D_2}\arrow[l,shift left=.5ex]\arrow[l,shift right=.5ex]\arrow[uu,bend right=16]
\end{tikzcd}}

Une telle donnée : \unsure{revient|équivaut} à la donnée
% At the foot only the tops of further strokes are visible.
$\unsure{F}\qquad\unsure{X}$
% archive sheet number: 15
